\section{Dialectics Redux}

\begin{quotex}
The premise from which the Buddhist Doctrine of Awakening starts is the destruction of the demon of dialectics; the renunciation of the various constructions of thought and speculation which are simply an expression of opinion, and of the profusion of theories, which are projections of a fundamental restlessness in which a mind that has not yet found itself its own principle seeks for support.
\flright{\textsc{Julius Evola}, \emph{The Doctrine of Awakening}}
\end{quotex}

\paragraph{Aporia}

Given the propensity for some readers to insist on the virtue and necessity of debate, despite our numerous posts to the contrary on Gornahoor, it may be illustrative to give an example. As we have mentioned, the Hermetic method is the way of depth, not confrontation. If we have two opposing view, each of which seems reasonable or even unrefutable, then we try to go deeper to determine what they have in common. A frequent example involves the debate between believers in religion and believers in science. When we penetrate to their unexamined assumptions, we find two fundamental questions which are at the core of their disagreement:

\begin{table}[h]\centering
\begin{tabular}{ccc}\toprule
~ & Religion & Science \\\midrule
What is Immutable & God & Matter \\
What is the cause of motion & Soul & Mechanical laws \\\bottomrule
\end{tabular}
\end{table}


Here we see two fundamental and incompatible views of the world which can never be resolved by debate on the level of thought.

For the believer in science, Matter (including both mass and energy) is the only necessary being, it is immutable since it cannot be created or destroyed and the total amount of Matter is constant. Everything that happens is the result of mechanical laws that can be discovered by the scientist. But how the immutable Matter (at the Big Bang) began to move is inexplicable.

For the religious believer, the immutable has to be outside of and transcendent to the world process, and this he calls God. When it comes to motion, he believes it to arise from intent, as he experiences every day. Thus he believes he has a soul that initiates motion that can't be explained from within the material world process.

It is obvious that at this level, both perspectives are reduced to the point of aporia so their argument is pointless and moot. Just as the believer can offer no explanation of God, neither can the scientist tell us what Matter is. While the believer cannot account for the motions of the planets (unless he resorts to angels and gods as the ancients did), the scientist cannot account for our experience of free will. They have no choice but to regard each other's viewpoint as delusional.

\paragraph{Gnosis}


So how do we resolve this? For the Hermetist, there is a higher form of knowledge, gnosis, about thought and sense experience, although it is more like the latter than the former. Just as I ``know'' there is a banana tree outside my window because I can perceive it directly, the Hermetist knows things of a metaphysical nature because he can perceive them directly, although not in the manner of sense perception, but certainly not by just thinking of them conceptually.

So, for example, by meditation practice, the Hermetist learns to observe the inner workings of consciousness. He may notice that while the contents of consciousness may be in motion or perpetual change, he also may notice that there is a constant, something residual that remains unchanged and unmoved, transcendent to the contents since it is not itself part of that content. Hence, the Hermetist knows by direct gnosis the immutable and that which is in motion. From that observation he builds and thinks.


\flrightit{Posted on 2011-09-11 by Cologero}

